\section{Por qué el Transformer es adecuado para los datos biomédicos}

La sección anterior estableció el mapa del curso; esta sección responde la pregunta común desde la capa de representación. La ventaja más directa del Transformer no es que «entienda medicina», sino que puede establecer dependencias contextuales en secuencias largas y recibir, con un esqueleto de cómputo similar, token discretos, eventos, aminoácidos o nucleótidos. Esta perspectiva unificada es poderosa, pero solo es válida si al mismo tiempo se conserva la semántica de cada dominio.

\subsection{Los cuatro tipos de secuencia no son el mismo lenguaje}

El ponente primero pide a los estudiantes presentes que respondan por qué el Transformer es adecuado para la biomedicine y luego aterriza las respuestas en cuatro objetos concretos. La clinical note es una secuencia de texto escrita por el médico; el electronic medical record (EMR, expediente clínico electrónico) puede verse como una secuencia de eventos formada por las interacciones del paciente con el sistema de salud; la protein es una secuencia de aminoácidos unidos por enlaces peptídicos; y el genome es una secuencia de bases nucleotídicas. Todas tienen un orden, pero poseen «vocabularios», escalas temporales y fuentes de error completamente distintos.

\lecturefigure{slide-03-why-transformers-biomedicine.jpg}{Pregunta de primeros principios: ¿por qué el Transformer es adecuado para la biomedicina?}{Diapositiva V003 recuperada del video oficial; Stanford CS25 Lecture 17}

\lecturefigure{slide-04-clinical-notes-sequences.jpg}{Notas clínicas: secuencias de texto formadas por el lenguaje de los médicos}{Diapositiva V004 recuperada del video oficial; Stanford CS25 Lecture 17}

\begin{knowledgebox}{Lectura de la figura: primero distinguir la «unidad de secuencia»}
En las notas clínicas, la posición suele indicar el orden del texto, pero en un registro de eventos la posición puede representar tiempo, consultas o estudios; en las proteínas, la posición corresponde a un amino acid residue, y en el genoma, a un nucleotide. Lo común es que el contexto cambia el significado de la posición actual; la diferencia es que el modelo debe usar tokenización, codificación posicional, función objetivo y formas de verificación adaptadas al dominio.
\end{knowledgebox}

La diapositiva del expediente clínico electrónico enfatiza «encounters with the medical system»: los diagnósticos, recetas, análisis de laboratorio y consultas del paciente no son solo párrafos de lenguaje natural, sino un flujo de eventos con tiempos, intervalos irregulares y valores faltantes. La diapositiva de proteínas desciende hasta el nivel molecular: la disposición lineal de los aminoácidos se pliega en una estructura tridimensional que influye en los sitios de unión y en la función; por eso, el modelado de secuencias puede aportar representaciones, pero no puede prescindir de la estructura ni de los experimentos.

\lecturefigure{slide-05-electronic-medical-records.jpg}{Expediente clínico electrónico: secuencia de eventos de las interacciones del paciente con el sistema de salud}{Diapositiva V005 recuperada del video oficial; Stanford CS25 Lecture 17}

\lecturefigure{slide-06-protein-sequences.jpg}{Proteínas: secuencias biológicas formadas por aminoácidos unidos}{Diapositiva V006 recuperada del video oficial; Stanford CS25 Lecture 17}

\teachervoice{El ponente llama en broma a las clinical notes «doctor gibberish» y enseguida se corrige a doctor speak. La broma encierra un problema real de diseño: quizá el modelo pueda generar terminología médica, pero si el paciente puede entenderla y si el médico puede verificarla con rapidez son dimensiones que la human evaluation posterior debe medir por separado.}

\subsection{El genoma y «las secuencias están en todas partes»}

La diapositiva del genoma muestra pares de bases nucleotídicas y, a continuación, reúne texto clínico, imágenes, proteínas y DNA en un mismo diagrama general. Aquí, multimodal (multimodal) no significa forzar todos los datos a convertirse en una oración, sino permitir que cada modalidad conserve su propia codificación y luego establecer relaciones mediante una representación común o atención entre modalidades. Por ejemplo, una imagen no es por naturaleza una secuencia unidimensional, pero puede dividirse en patch token; las trayectorias de expresión génica tampoco son lenguaje, pero pueden servir como objetivo de predicción en cada posición.

\lecturefigure{slide-07-genome-sequences.jpg}{Genoma: secuencias largas formadas por pares de bases nucleotídicas}{Diapositiva V007 recuperada del video oficial; Stanford CS25 Lecture 17}

\lecturefigure{slide-08-biomedical-sequences-everywhere.jpg}{En la biomedicina, las secuencias y las modalidades están en todas partes}{Diapositiva V008 recuperada del video oficial; Stanford CS25 Lecture 17}

\begin{longtable}{@{}L{0.16\textwidth}L{0.19\textwidth}L{0.22\textwidth}L{0.31\textwidth}@{}}
\toprule
Término & Unidad de secuencia & Contexto principal & Restricciones del dominio que no se pueden ignorar \\
\midrule
clinical note & subword / word & Antecedentes y problema actual & Abreviaturas, negaciones, tiempo, privacidad, consenso clínico \\
EMR event & Consulta/diagnóstico/fármaco/análisis de laboratorio & Trayectoria longitudinal del paciente & Tiempos irregulares, faltantes, sistemas de codificación y sesgo de selección \\
protein sequence & amino acid & motif locales y relaciones entre residuos distantes & Plegamiento tridimensional, conservación evolutiva, función experimental \\
genome sequence & nucleotide & Elementos reguladores e interacciones de largo alcance & Regiones no codificantes, tipo celular, entorno de la cromatina y ruido de medición \\
multimodal input & text/image/signal token & Relaciones dentro de cada modalidad y entre modalidades & Alineación, sesgo de muestreo, escalas de error distintas \\
\bottomrule
\end{longtable}

Este glosario implica además una conclusión importante de ingeniería de datos: el llamado «preentrenamiento a gran escala» no es lo mismo en los cuatro tipos de datos. La escala del texto clínico está limitada por la privacidad, los hábitos de redacción de cada institución y la desidentificación; los eventos del EMR suelen presentar sesgo de selección por los seguros, el acceso a la atención y los sistemas de codificación; muchas secuencias de las bases de datos de proteínas provienen de especies cercanas, por lo que una partición aleatoria puede hacer que el conjunto de entrenamiento y el de prueba compartan parientes cercanos; y los track genómicos se concentran en unas pocas líneas celulares y protocolos experimentales. Si solo se informa el número de token y no la procedencia, la deduplicación o la partición por familia o por individuo, el modelo puede obtener resultados optimistas al memorizar muestras similares. Por lo tanto, el primer trabajo de ingeniería después de la representación de secuencias debe ser definir la independencia: qué fuga bloquea cada partición, ya sea por paciente, por tiempo, por familia de proteínas o por cromosoma.

\subsection{Esqueleto de cómputo común y límites de aplicabilidad}

La diapositiva general condensa las ventajas en tres términos: sequence modeling, long-range dependencies, scaling. La self-attention estándar construye similitudes por pares para una secuencia de longitud $n$, por lo que puede conectar directamente posiciones lejanas; además, el preentrenamiento permite al modelo aprovechar grandes cantidades de datos sin etiquetar o con etiquetas débiles. Pero las ventajas computacionales y estadísticas no traen automáticamente seguridad clínica, veracidad de la función proteica ni causalidad de las variantes; los cuatro casos que siguen buscan precisamente cubrir estas carencias.

\lecturefigure{slide-09-transformer-biomedical-fit.jpg}{Puntos de coincidencia entre el Transformer y los datos biomédicos}{Diapositiva V009 recuperada del video oficial; Stanford CS25 Lecture 17}

\begin{importantbox}{Representación unificada no equivale a evidencia unificada}
Lo que el Transformer aporta es un lenguaje de cómputo unificado: representación por token, interacción contextual y entrenamiento escalable. Las conclusiones médicas siguen requiriendo un ciclo cerrado de evidencia del dominio: las respuestas clínicas las evalúan médicos y pacientes, las descripciones de proteínas deben contrastarse con bases de datos y experimentos, la corrección de errores de secuenciación debe medirse con read-level accuracy, y las predicciones regulatorias deben compararse con genomic tracks reales y variant assays.
\end{importantbox}

\begin{warningbox}{El error de pensar que «todo es una secuencia»}
Colocar los objetos en una fila no significa que el modelo haya aprendido los invariantes correctos. Los intervalos de tiempo del EMR, la proximidad tridimensional de las proteínas y la dependencia del tipo celular en el DNA pueden quedar ocultos por un orden unidimensional. La serialización es el punto de entrada del modelado, no un sustituto de la estructura del dominio.
\end{warningbox}

\subsection{Resumen de la sección}

El texto clínico, los eventos médicos, las proteínas y el genoma comparten el orden y la dependencia del contexto, pero difieren por completo en la semántica de los token, la longitud, el ruido y la forma de validación. El valor del Transformer está en ofrecer un esqueleto de cómputo transferible; un sistema biomédico de alta calidad, en cambio, debe restituir para cada tipo de datos la estructura del dominio y la evidencia independiente.
